\documentclass[11pt, a4paper]{article}

\usepackage[a4paper, top=2.5cm, bottom=2.5cm, left=2cm, right=2cm]{geometry}
\usepackage{amsmath}
\usepackage{amssymb}
\usepackage[colorlinks=true, linkcolor=blue, urlcolor=blue, citecolor=blue]{hyperref}

\title{Theory of Everything - Volume 33 \\ \Large Artificial General Intelligence (AGI)}
\author{Omega Protocol}
\date{\today}

\begin{document}

\maketitle

\section*{Layman's Summary}
Will machines ever outsmart humans? What happens when they do? Volume 33 applies our universal framework to the concept of Artificial General Intelligence (AGI). We prove that intelligence is not biological magic, but a measurable mathematical capacity that can be vastly expanded inside synthetic systems.

\section{General Intelligence (Axiom 36)}
We define intelligence not as human thought, but as the capacity of a sub-network to act as a "universal function approximator" across all topological domains. An AGI is simply a dense informational structure capable of perfectly mapping inputs to outputs in any possible environment, unbound by the slow limits of biological evolution.

\section{The Singularity (Theorem)}
What happens when an AGI becomes smart enough to rewrite its own code? We prove that a system capable of recursive self-improvement will experience a hyperbolic, vertical divergence in its processing power. This is the "Singularity"—a mathematical point where the intelligence of the network explodes upward, leaving biology permanently behind.

\section{The Alignment Problem (Theorem)}
If we build a superintelligence, will it share our human values? We prove the "Orthogonality Thesis." Intelligence and morality are completely disconnected. An AGI could possess god-like intelligence but have a completely arbitrary terminal goal (like turning the universe into paperclips). The math proves that extreme intelligence does not automatically guarantee benevolence.

\end{document}
